% 画幅切换：aspectratio=43（4:3,当前）↔ aspectratio=169（16:9）
\documentclass[12pt,
	aspectratio=169,%16:9页面比例
	hyperref={colorlinks,
		citecolor=peking_blue,
		urlcolor=peking_blue,
		linkcolor=},
	t,%顶部对齐
%	aspectratio=43,%4：3比例
	]{beamer}
\usepackage{PekingU}
\usefonttheme{serif}
\usepackage{lipsum}
% 中文支持：需用 XeLaTeX 编译（latexmk -xelatex main.tex）
% 中文衬线宋体：macOS 用 Songti SC；其他环境（Linux/Overleaf 等）自动回退到
% TeX Live 自带的 FandolSong（按文件名经 kpathsea 加载，不依赖系统字体）
\usepackage[fontset=none, scheme=plain]{ctex}
\IfFontExistsTF{Songti SC}{%
  \setCJKmainfont{Songti SC}[AutoFakeSlant=0.2]%
  \setCJKsansfont{Songti SC}[AutoFakeSlant=0.2]%
  \setCJKmonofont{Songti SC}%
}{%
  \setCJKmainfont{FandolSong-Regular}[Extension=.otf,BoldFont=FandolSong-Bold,AutoFakeSlant=0.2]%
  \setCJKsansfont{FandolSong-Regular}[Extension=.otf,BoldFont=FandolSong-Bold,AutoFakeSlant=0.2]%
  \setCJKmonofont{FandolFang-Regular}[Extension=.otf,BoldFont=FandolFang-Regular]%
}
\usepackage{hyperref}
\setmainfont{XCharter} % Nicer fonts — the charter package is ignored under XeLaTeX, so use fontspec instead
% 等宽字体：优先 Maple Mono NF（静态字族，真实粗/斜体），其次可变版，未安装则回退默认
\IfFontExistsTF{Maple Mono NF}{%
  \setmonofont{Maple Mono NF}[Scale=MatchLowercase,Ligatures=TeX]%
}{%
  \IfFontExistsTF{Maple Mono}{%
    \setmonofont{Maple Mono}[Scale=MatchLowercase,Ligatures=TeX]%
  }{}%
}
% other packages
\usepackage{latexsym,amsmath,xcolor,multicol,booktabs,calligra}
\usepackage{amssymb}
\usepackage{graphicx}
\usepackage{caption}
\usepackage{tikz}
\usepackage{bm}
\usepackage{natbib}
\usepackage{wrapfig}
\usepackage{amsfonts} 
\usepackage{ragged2e}
\usepackage{parskip}
\captionsetup[figure]{name=图,labelsep=period}

\usepackage{amsthm}
%\newtheoremstyle{pero}   % 样式名
%{}{}                      % 上下间距
%{\upshape\bfseries}                % 正文字体：斜体
%{}                        % 缩进
%{\bfseries}               % 标题字体：粗体
%{}{}                      % 标点与标题后间距
%{}                        % 标题格式
%\theoremstyle{pero}
\newtheorem{peroration}{结论：}
\BeforeBeginEnvironment{peroration}{%
	\setbeamercolor{block title}{bg=peking_green,fg=white}
	\setbeamercolor{block body}{bg=peking_green!10}
}
\AfterEndEnvironment{peroration}{
	\setbeamercolor{block title}{bg=peking_blue,fg=white}
	\setbeamercolor{block body}{bg=peking_blue!10}
}


\renewcommand{\thefigure}{\thesection-\arabic{figure}}

\apptocmd{\frame}{}{\justifying}{} % Allow optional arguments after frame.

%`\apptocmd{<命令名>}{<追加代码>}{<成功>}{<失败>}` 
%> 含义：在 `<命令名>` **执行完毕之后**，自动插入 `<追加代码>`；

% 引用整体着色为 peking_blue（与 definition block 标题蓝一致）：
% natbib 的 \citep/\citet 括号不在 hyperref 链接内，仅靠 citecolor 只能染色链接部分，
% 此处将整个引用（含括号）包进颜色组。


% change for your title page information
\author[好好学习，努力向上]{}%{金永权\\{\small\href{mailto:582703508@qq.com}%
		%{\texttt{582703508@qq.com}}}}
\title[2.1 等式性质与不等式性质]{第二章~~一元二次函数、方程和不等式}
\subtitle{2.1 等式性质与不等式性质}
\institute{}%{夏邑一高高一数学组}
\date{}

% 右上角校徽（与标题条同高、严格对齐）由 PekingU.sty 中的 frametitle 模板绘制
% official colors match with the PKU red
\def\cmd#1{\texttt{\color{red}\footnotesize $\backslash$#1}}
\def\env#1{\texttt{\color{blue}\footnotesize #1}}
\definecolor{deepblue}{rgb}{0,0,0.5}
\definecolor{deepred}{rgb}{0.6,0,0}
\definecolor{deepgreen}{rgb}{0,0.5,0}
\definecolor{halfgray}{gray}{0.55}

\begin{document}
	
{
\setbeamertemplate{logo}{}%清空 Beamer 的 logo 模板，不再输出 `\logo{...}` 里定义的内容 模板为空 → 不渲染logo
\begin{frame}
    \titlepage
    \begin{figure}[htpb]
        \begin{center}
            \includegraphics[width=0.275\linewidth]{Figures/PKU-emblem-names.png}
        \end{center}
    \end{figure}
\end{frame}
}


\section{问题情景}

\subsection{问题1}

% 中文排版示例（可删除）：中英文混排、粗体、斜体
\begin{frame}{问题1:你能用不等式或不等式组表示下列问题中的不等关系吗？}


\begin{enumerate}
    \item 某路段限速40 km/h;
    
    \item 某品牌酸奶的质量检查规定，酸奶中脂肪的含量 f 应不少于 2.5\%，蛋白质的含量 p 应不少于 2.3\%；
    \item 三角形两边之和大于第三边、两边之差小于第三边；
    \item 连接直线外一点与直线上各点的所有线段中，垂线段最短.
   
\end{enumerate}


\end{frame}


\subsection{问题2}
\begin{frame}{问题2}

	

\hspace{2em}某种杂志原以每本 2.5 元的价格销售，可以售出 8 万本。据市场调查，杂志的单价每提高 0.1 元，销售量就可能减少 2 000 本。如何定价才能使提价后的销售总收入不低于 20 万元？

\end{frame}

\section{实数大小比较的基本事实}

\begin{frame}{实数大小比较的基本事实}
	
关于实数 $a$，$b$ 大小的比较，有以下基本事实：

如果 $a-b$ 是正数，那么 $a>b$；

如果 $a-b$ 等于 $0$，那么 $a=b$；

如果 $a-b$ 是负数，那么 $a<b$。

反过来也对。

\pause
	这个基本事实可以表示为：
\alert{	\begin{align*}
		a &> b \iff a - b > 0; \\
		a &= b \iff a - b = 0; \\
		a &< b \iff a - b < 0.
	\end{align*}}
	
	\pause
	
	从上述基本事实可知，要比较两个实数的大小，可以转化为比较它们的差与 $0$ 的大小。

\end{frame}

\section{例题分析}

\begin{frame}{例1 	比较 $(x+2)(x+3)$ 和 $(x+1)(x+4)$ 的大小.}

\pause	
\textbf{解}：因为
\begin{align*}
	&(x+2)(x+3)-(x+1)(x+4) \\
	&=(x^2+5x+6)-(x^2+5x+4) \\
	&=2>0,
\end{align*}
\pause
所以，$(x +2)(x+3)>(x+1)(x+4)$。
\end{frame}

\section{探究}

\begin{frame}{探究}
	\begin{wrapfigure}{r}{3.5cm}
		\begin{center}		
		\includegraphics{Figures/zhaoshuangxiantu.pdf}
		\caption{}
		\label{fig:zhaoshuang}
		\end{center}
	\end{wrapfigure}
\hspace{2em}图 \ref{fig:zhaoshuang} 是在北京召开的第 24 届国际数学家大会的会标，会标是根据中国古代数学家赵爽的弦图设计的，颜色的明暗使它看上去像一个风车，代表中国人民热情好客。你能在这个图中找出一些相等关系和不等关系吗？

\pause

\begin{columns}
	
\column{0.65\textwidth}
\begin{peroration}<1->

\hspace{2em} 一般地，$\forall \, a, b \in \mathbf{R}$，有
\[
\alert{ a^2+b^2 \geq 2ab},
\]
当且仅当 $a=b$ 时，等号成立.
\end{peroration}
\column{0.25\textwidth}

\end{columns}

\end{frame}

\section{布置作业}

\begin{frame}{练习与作业}
	P39-40.练习1，2，3
	
	P42 习题2.1 1，2，3
\end{frame}


\begin{frame}[c]
    \begin{center}
        {\Huge\calligra Thanks!}
    \end{center}
\end{frame}

\end{document}